\section{Related Work}\label{sec:related}
\textbf{LLM trading agents.} FinMem~\cite{yu2025finmem} introduced layered memory and character design, FinCon~\cite{yu2024fincon} and TradingAgents~\cite{xiao2025tradingagents} coordinate multiple LLM roles, FinAgent~\cite{zhang2024finagent} adds multimodal tools, and InvestorBench~\cite{li2025investorbench} standardizes the perception--profile--memory--action design we follow. Finance-specific language models~\cite{araci2019finbert,wu2023bloomberggpt} and retrieval-augmented generation~\cite{lewis2020rag} supply the text channel these agents rely on.

\textbf{Attacks on LLM agents.} Knowledge-corruption attacks inject crafted passages into a retrieval corpus~\cite{zou2025poisonedrag}, and indirect prompt injection hides instructions in retrieved content~\cite{greshake2023indirect}. Automated red-teaming systems search over attack strategies~\cite{zhou2025autoredteamer,jiang2024artprompt}. For trading agents, TradeTrap~\cite{yan2025tradetrap} shows that small textual changes shift decisions, and AutoRedTrader~\cite{liu2026autoredtrader}, the attack we defend against, adds a finance-structured operator space with closed-loop feedback.

\textbf{Defenses.} Detection-based defenses try to flag manipulated text, but reference-free detection of financial misinformation remains hard~\cite{jiang2026glisters,wang2026life}, and an adaptive attacker can filter its candidates against the detector. Robust aggregation bounds the influence of a few corrupted passages, e.g.\ by isolating passages and aggregating their answers~\cite{xiang2024robustrag}, but assumes that most retrieved passages are clean. AutoRedTrader~\cite{liu2026autoredtrader} proposes time-series-informed grounding with 30-day windows and analog outcomes as a preliminary defense. We turn it into a complete layer with leak-free analogs, microstructure evidence and a risk layer, and evaluate each component under adaptive and replayed attacks. Position sizing and stop rules are standard risk controls in trading systems~\cite{lopezdeprado2018advances}; we test whether they also limit misinformation.

\textbf{Financial misinformation and methodology.} Media tone moves markets~\cite{tetlock2007giving} and dictionary methods quantify it~\cite{loughran2011liability}. Behavioral finance supplies the biases attackers exploit: overconfidence~\cite{barber2001boys}, loss aversion~\cite{kahneman1979prospect}, herding~\cite{spyrou2013herding} and noise trading~\cite{delong1990noise}. Negative controls detect confounding in observational studies~\cite{lipsitch2010negative}, which motivates our placebos. Lookahead and backtest overfitting are the classic failure modes of trading research~\cite{lopezdeprado2018advances,bailey2014probability}.
